\subsubsection{The Spectra of Bounded Operators}\label{subsec:bounded-operator-spectra}

We now study the properties of the spectra of bounded operators. Many of our most important results about the spectrum (including the results for the results below for compact operators) are derived by using a power series expansion for the resolvent. We now prove a fundamental theorem that is the analogue of the elementary calculus result on the convergence of geometric series.

\begin{Theorem}
\label{thm:neumann-operator-series}
Let $X$ be a Banach space and suppose $A\in\mathcal{L}(X)$ satisfies $\|A\|<1$. Then $(I-A)^{-1}$ exists and is bounded, and the following power series expansion for $(I-A)^{-1}$ converges in the operator norm:
\begin{equation}\label{eq:neumann-operator-series}
(I-A)^{-1}=\sum_{k=0}^\infty A^k=I+A+A^2+\cdots.
\end{equation}

\end{Theorem}

\begingroup
\postdisplaypenalty=10000
\renewcommand{\endProof}{\par\nobreak\hfill$\blacksquare$\par\vspace{0.1cm}}
\begin{Proof}
The main idea in this proof is that if a series in a Banach space \emph{converges absolutely} (i.e., the sum of the norms of the terms converges), then the original series converges. (The proof of this fact is identical to the elementary calculus proof for series of real numbers.) In our case, the Banach space in question is $\mathcal{L}(X)$, and we have
\begin{equation}\label{eq:neumann-series-norm-majorant}
\sum_{k=0}^\infty\|A^k\|\leq\sum_{k=0}^\infty\|A\|^k.
\end{equation}
Since $\|A\|<1$, the geometric series on the right converges. Hence, the series on the right of \eqref{eq:neumann-operator-series} is absolutely convergent and therefore convergent. We need only show that its limit is indeed $(I-A)^{-1}$, and the proof of this is essentially the same as the elementary calculus result for geometric series; i.e., we have\[\begin{aligned}
I-A^{k+1}&=(I-A)(I+A+A^2+\cdots+A^k)\\
&=(I+A+A^2+\cdots+A^k)(I-A).
\end{aligned}\]Now, since $\|A\|<1$ we have $\lim_{k\to\infty}A^{k+1}=0$. Thus,
\begin{equation}\label{eq:neumann-series-two-sided-inverse}
I=(I-A)\left(\sum_{k=0}^\infty A^k\right)=\left(\sum_{k=0}^\infty A^k\right)(I-A),
\end{equation}
and the theorem is proved.
\end{Proof}
\endgroup

This theorem immediately gives us the following result, which says that the spectrum $\sigma(A)$ of a bounded operator $A$ lies in a bounded disk in the complex plane.

\begin{Corollary}
\label{cor:bounded-operator-spectrum-norm-bound}
Let $A\in\mathcal{L}(X)$, and suppose $\lambda\in\sigma(A)\subset\mathbb{C}$. Then
\begin{equation}\label{eq:bounded-operator-spectrum-norm-bound}
|\lambda|\leq\|A\|.
\end{equation}

\end{Corollary}

\begingroup
\postdisplaypenalty=10000
\renewcommand{\endProof}{\par\nobreak\hfill$\blacksquare$\par\vspace{0.1cm}}
\begin{Proof}
Suppose $|\lambda|>\|A\|$. Then we can show that $\lambda\in\rho(A)$ by using Theorem~\ref{thm:neumann-operator-series} to construct the resolvent as follows:
\begin{equation}\label{eq:resolvent-series-outside-norm-disk}
R_\lambda(A)=(A-\lambda I)^{-1}=-\dfrac{1}{\lambda}\left(I-\dfrac{1}{\lambda}A\right)^{-1}=-\dfrac{1}{\lambda}\sum_{k=0}^\infty\left(\dfrac{1}{\lambda}A\right)^k.
\end{equation}
Here we have used the fact that $\left\|\dfrac{1}{\lambda}A\right\|<1$. This completes the proof.
\end{Proof}
\endgroup

Since we have just shown that the spectrum of a bounded operator is contained in a disk, it is natural to ask whether this disk is optimal. Thus, we give the following definition.

\begin{Definition}
\label{def:operator-spectral-radius}
The \textbf{spectral radius} of an operator from $X$ to $X$ is defined to be
\begin{equation}\label{eq:operator-spectral-radius-definition}
r_\sigma(A):=\sup_{\lambda\in\sigma(A)}|\lambda|.
\end{equation}

\end{Definition}

Thus, for $A\in\mathcal{L}(X)$, Corollary~\ref{cor:bounded-operator-spectrum-norm-bound} simply says
\begin{equation}\label{eq:operator-spectral-radius-norm-bound}
r_\sigma(A)\leq\|A\|.
\end{equation}
In general, equality does not hold in \eqref{eq:operator-spectral-radius-norm-bound}, but it does hold for a class of operators called \emph{normal}. Problem~\ref{ex:self-adjoint-spectral-radius-equals-norm} below establishes equality for self-adjoint operators.
